%------------------%
% 10 gener 2020    %
%------------------%
 Marina, per posar un títol com el que té el document, no cal fer servir \maketitle. Això es fa servir per fer la portada d'un article, llibre, etc. (En el cas que ens ocupa n'hi ha prou de jugar amb el tamany del text, centrar-lo i ja fa de títol). Jo personalment, per fer arxius petits, estil exàmens, etc, faig servir la classe "article" perquè ja em va bé les mides de marges que posa, etc. Pots explorar si vols les diferents opcions que té el "documentclass". Cadascuna imposa un format general característic a les pàgines del document. Per exemple, la classe "book", fa servir marges laterals diferents a les pàgines senars i parells. Això es fa perquè quan llegim un llibre, les pàgines "de la dreta" han de tenir prou marge a l'esquerra perquè el text no quedi tapat pel plec del llibre. Com veus, TeX contempla normes molt rigoroses i detallistes sobre el "layout" dels documents, i és dels pocs softwares gratuïts que respecte les normes tipogràfiques a l'hora d'escriure. Això fa que el resultat tingui aquest aspecte professional i que es faci servir a nivell mundial en les comunitats universitàries científiques.

Fixa't que el compilador et marca un "warning" (color taronja). Això vol dir que hi ha alguna errada al codi. Com és un "warning" i no un "error", el compilador produeix un resultat. Segons la gravetat dels errors, a vegades no surt res, o es veuen coses estranyes al document.

Si fas clic a l'advertiment, sovint et donen pistes sobre el que està passant, la línea de codi que pot contenir l'error, etc. Mira-ho i si no ho veus clar, m´ho dius.
Una altra cosa és el tema de l'idioma. Faràs el Tr en castellà, en català? Tria perquè segons això hauràs de fer servir un "babel" o un altre. És important a l'hora de posar accents, etc. 

Acaba de fer l'examen aquest de mostra, revisa si vols el "lshort" que tens a la carpeta del drive i mira de reproduir l'altre article. Per fer-lo hauràs d'experimentar amb moltes coses més, a banda d'un munt de símbols. T'anirà bé per practicar amb Tex. Pot ser una mica feixuc al principi, però els resultats valen la pena. Ves practicant amb Tex, però també ves mirant el document sobre criptografia... Al final hauràs de fer la primera part de la memoria a partir de la informació i exemples que et proporciona el document, i els exercicis del final de l'arxiu seran part del treball de camp. Més endevant, començarem a mirar lo de les corbes el·líptiques.

%-------------------------------------------------%
